\section{An original-law reduction for four summands}
\label{r4bridge:section}

This appendix proves a conditional analytic implication for the finite base.
It does not supply a complete four-summand certificate. The smoothing
inequality used below is proved in Appendix~\ref{v9:prawitz:appendix}. The remaining steps are proved for the original iid marginal, including laws with
infinite support inside a bounded interval and atoms at zero or at a
threshold. No optimizer or contact assumption is used.

Let $X$ have mean zero, variance one, $|X|\le R=17/4$, and
$\rho=\E|X|^3<46/25$. Set
\[
 u=\E X_+=\E X_->0,\quad v_\pm=\E X_\pm^2,
 \quad a=v_-/u,\quad b=v_+/u,
\]
\[
 w=u/a+u/b,\qquad \vartheta=v_+-v_-.
\]
Both $v_\pm$ and $a,b$ are positive; also $a,b\le R$.
Signwise Cauchy gives $u/a\le\Pp(X<0)$ and
$u/b\le\Pp(X>0)$, so $0<w\le1$. The canonical probability
\[
 Q=(u/a)\delta_{-a}+(1-w)\delta_0+(u/b)\delta_b
\]
has the same first and second moments on each strict half-line as $P$.
Put $r=\rho(Q)$ and $\eta=\rho-r$.

\subsection{One actual measure, including all contact atoms}
For $y=|X|$, let $c=b$ on the positive tag and $c=a$ on the
negative tag, and put $W=y(y-c)^2$ (zero at $X=0$). Direct expansion gives
\begin{align}
 \E W&=\eta\ge0,\label{r4bridge:excess}\\
 \E X^3&=b-a+\E[\operatorname{sgn}(X)W],\notag\\
 \E X^4&=a^2-ab+b^2+\E[(y+2c)W].\label{r4bridge:moments}
\end{align}
Indeed $y^3=W+2cy^2-c^2y$ and
$y^4=(y+2c)W+3c^2y^2-2c^3y$; integrate on each tag using
$\E[y;\text{tag}]=u$ and $\E[y^2;\text{tag}]=uc$.
If $\eta=0$, then $P$ is supported on $\{-a,0,b\}$ and its
first two moments force $P=Q$. This case requires no defect probability.

For $\eta>0$, define the single tagged probability $\nu=W\,dP/\eta$.
It charges neither zero nor either contact. On its domain set
\[
 k_0=(c^2y)^{-1},\qquad k_2=(y-c)^{-2},\qquad k_1=(y-c)^{-1}.
\]
Then, writing $p_0=1-w$,
\begin{align}
 \eta\int k_0\,d\nu&=p_0-P\{0\},\notag\\
 \eta\int_{\sigma} k_2\,d\nu&=u-c_\sigma P\{\sigma c_\sigma\},
 &\int_{\sigma}k_1\,d\nu&=0.\label{r4bridge:budgets}
\end{align}
To check the first formula, $Wk_0=(y-c)^2/c^2$ off zero and
contacts; its integral on a tag is $P(\text{tag})-u/c$.
For $k_2$, cancellation gives $y$ only away from the contact,
so the omitted contact mass must be subtracted. For $k_1$ it gives
$y(y-c)$, whose integral is zero even after omitting the contact.
All integrals exist; for $k_1$, Cauchy uses the finite $k_2$ integral.

Let $\psi(q)=\E e^{\mathrm iqX}$ and $g(q)=\int e^{\mathrm iqx}\,dQ(x)$.
Hermite interpolation of $f(y)=e^{\mathrm i\sigma qy}$ at $0,c,c$ gives
\[
 f(y)-L_f(y)=y(y-c)^2 r_q(\sigma,y),\qquad
 r_q=f[0,c,c,y].
\]
The interpolant is quadratic with constant $f(0)=1$. Its linear and
quadratic integrals match on each tag; its constants match only after
both tags and the zero atoms are combined. Consequently
\begin{equation}
 \psi(q)-g(q)=\eta\int r_q\,d\nu.\label{r4bridge:response}
\end{equation}
This identity uses the same measure at all frequencies. The integral
form of a third divided difference gives $|r_q|\le |q|^3/6$.
For completeness another bound is $|r_q|\le q^2/y$: write
\[
 L_f(y)=f(c)+f'(c)(y-c)
       +\frac{f(0)-f(c)+cf'(c)}{c^2}(y-c)^2.
\]
The quadratic coefficient and the first-order Taylor remainder at $c$
are bounded by $q^2/2$ and $q^2(y-c)^2/2$, respectively. Division
gives the claim. Thus the following three radii are valid:
\begin{equation}
 |\psi(q)-g(q)|\le
 \min\{\eta|q|^3/6,\ \max(a,b)^2p_0q^2,\ 1+|g(q)|\}.
 \label{r4bridge:radii}
\end{equation}
The mass radius follows from $1/y=c^2k_0$ and \eqref{r4bridge:budgets}.

For arbitrary finite coefficients $\kappa_j$, define
$A=\Re\sum_j4\kappa_jg(q_j)^3r_{q_j}$. If fixed prices
$\lambda_0,\lambda_+,\lambda_-\ge0$ and signed $\theta_\pm$ satisfy
\[
 A-\lambda_0k_0-\lambda_\sigma k_2-\theta_\sigma k_1\le S
\]
on the entire tagged domain, then
\begin{equation}
 \Re\sum_j4\kappa_jg(q_j)^3(\psi(q_j)-g(q_j))
 \le p_0\lambda_0+u(\lambda_++\lambda_-)+S\eta=:\Gamma(\eta).
 \label{r4bridge:price}
\end{equation}
This holds for each actual $P$, not for a purported sufficiency
characterization of all measures satisfying the budgets. Enlarging
that class in an upper-bound problem is safe. Price choices must precede
the supremum; signed equality prices and a negative $S$ are retained.
At $\eta=0$ the response is zero and the displayed right side is
nonnegative, so the same bound applies directly.

\subsection{The physical smoothing and its two reserves}
Let $Z=(X_1+\cdots+X_4)/2$ with independent copies of this same law.
The input Prawitz inequality is, for every real $x$,
\begin{equation}
 \Pp(Z\le x)\le J_P(x):=\frac12+
 \int_0^1\Re\{\widehat K(t)e^{-6\mathrm itx}\psi(3t)^4\}\,dt,
 \quad \widehat K(t)=1-t+\mathrm i[(1-t)\cot(\pi t)+1/\pi].
 \label{r4bridge:prawitz}
\end{equation}
Only the real integrand is integrated. Put $S(t)=t\Im\widehat K(t)$.
Its endpoint-continuous extension satisfies
\begin{equation}
 0\le S\le1/3,\qquad |S'|<2,\qquad |S''|<5.
 \label{r4bridge:kernel}
\end{equation}
Here are explicit sufficient bounds. The cotangent partial fractions give
\[
 S=\pi^{-1}-2\pi^{-1}(A+MG),\quad
 A=t^2/(1+t),\ M=t^2(1-t),\
 G=\sum_{j\ge2}(j^2-t^2)^{-1}.
\]
On $[0,1]$,
$|A'|\le3/4$, $|A''|\le2$, $|M|\le1/6$,
$|M'|\le1$, $|M''|\le4$, and
$G\le3/4$, $|G'|\le10/27$, $|G''|\le106/135$.
The first sum telescopes; the latter two follow from
$j^2-t^2\ge3j^2/4$ and
$\sum_{j\ge2}j^{-k}\le2^{-k}+\int_2^\infty s^{-k}\,ds$.
The derivative series converge uniformly. With $\pi>3$, these imply
$|S'|<19/18$, $|S''|<223/54$. Nonnegativity follows directly from
$(1-t)\cot(\pi(1-t))\le1/\pi$ when $t>1/2$, and from the
nonnegative cotangent when $t\le1/2$. The positive partial-fraction
terms give $S\le1/\pi<1/3$.

Set $D=x-Z$, $A_x(t)=\E\cos(6Dt)$ and
$L_x(t)=\E[6D\operatorname{sinc}(6Dt)]$. The regular real integrand
is $(1-t)A_x(t)+S(t)L_x(t)$. Since
$\operatorname{sinc}(v)=\int_0^1\cos(sv)\,ds$,
\[
 |A_x'|\le6m_1,\ |A_x''|\le36m_2,\quad
 |L_x|\le6m_1,\ |L_x'|\le18m_2,\ |L_x''|\le72m_3,
 \quad m_j=\E|D|^j.
\]
Differentiation under these bounded-moment integrals and the product rule
therefore give
\begin{equation}
 |F_x''(t)|\le42m_1+108m_2+24m_3.\label{r4bridge:F2}
\end{equation}
All four copies are used to obtain
\[
 \E D^2=x^2+1=:D_2,\qquad
 \E D^4=x^4+6x^2-2x\E X^3+(\E X^4+9)/4.
\]
Thus for valid whole-cell moment bounds $\kappa_*\ge|\E X^3|$ and
$M_*\ge\E X^4$, put
$D_4=x^4+6x^2+2|x|\kappa_*+(M_*+9)/4$.
Then $m_1\le\sqrt{D_2}$ and $m_3\le\sqrt{D_2D_4}$.
The bounds \eqref{r4bridge:moments} supply these without an additional
fourth-moment premise. On $|x|\le2$, $m_2\le5$, $\E D^4<54$,
$m_1<9/4$, $m_3<33/2$, so \eqref{r4bridge:F2} is less than $2061/2$.

Choose $N\ge1$, $t_j=(j+1/2)/N$, $q_j=3t_j$ and
$\kappa_j(H)=\widehat K(t_j)e^{-\mathrm iq_jH}/N$.
Composite midpoint quadrature has error at most $\sup|F''|/(24N^2)$.
Consequently the raw numerator $2(J_P(H/2)-\Phi(H/2))$ has the
single reserve
\begin{equation}
 Q_H=\min\left\{\frac{687}{8N^2},
 \frac{42\sqrt{D_2}+108D_2+24\sqrt{D_2D_4}}{12N^2}\right\}.
 \label{r4bridge:quad}
\end{equation}
This quadrature is for the original law once, before reference expansion.
Neither the reference nor its response/remainder needs another quadrature.

To control threshold joins, suppose $a_j\ge|g(q_j)|$ and $e_j$
bounds $|\psi(q_j)-g(q_j)|$ uniformly on the entire excess cell.
The actual characteristic function is $1$-Lipschitz, since
$\E|X|\le1$. Hence
\[
 P_j=\min(1,a_j+e_j+3/(2N))
\]
bounds $|\psi(3t)|$ throughout its frequency cell. Define $d_j$ as
the smaller of the two integrals over that cell of
\[
 t^2(1-t)+t/3,\qquad t^2(1-t)+t(1-t)/2.
\]
Both dominate $t^2|\widehat K(t)|$: use
$|\widehat K|\le1-t+\Im\widehat K$ and, for the second,
$S\le(1-t)(1+2t)/(\pi(1+t))\le(1-t)/2$.
Their full integrals are $1/4$ and $1/6$. Thus
\[
 |J_P''(x)|\le B:=36\sum_jd_jP_j^4.
\]
For the smooth physical function $G_P(H)=2(J_P(H/2)-\Phi(H/2))$,
$|G_P''|\le(B+1/4)/2$, since $\sup|\Phi''|<1/4$.
On any raw-threshold interval of length $h$, its chord lies above
$G_P$ up to the single reserve
\begin{equation}
 E_H=h^2(B+1/4)/16.\label{r4bridge:H}
\end{equation}
There is no differentiation or interpolation of the CDF. The moment
reserve \eqref{r4bridge:quad} may be computed separately at each node;
the curvature reserve is common to the whole geometry/excess cell.

\subsection{A precise conditional certificate implication}
\begin{theorem}[Low-excess original-law bridge]
Fix a geometry cell, an excess interval $[L,U]\subset[0,\infty)$,
and raw nodes $-4\le H_0<\cdots<H_m\le25/8$ of spacing at most $h$.
Assume all the following bounds are valid for every actual standardized
$P$ in this cell, with $\rho=r+\eta<46/25$ and $\eta\in[L,U]$.
At each node let $M_H\ge C_H(Q)-Tr$, where
\[
 C_H(Q)=1+2\Re\sum_j\kappa_j(H)g(q_j)^4-2\Phi(H/2).
\]
Let $\Gamma_H(\eta)$ be a whole-domain price bound
\eqref{r4bridge:price}, and let $R_{2,H}(\eta)$ bound
\[
 \sum_j|\kappa_j(H)|
 |\psi(q_j)^4-g(q_j)^4-4g(q_j)^3(\psi(q_j)-g(q_j))|.
\]
Suppose the complete row
\[
 P_H(\eta)=M_H+Q_H+E_H+2\Gamma_H(\eta)+2R_{2,H}(\eta)-T\eta
\]
is a convex quadratic, nonpositive at $L$ and $U$, at every node.
Then $2\{\Pp(X_1+\cdots+X_4\le H)-\Phi(H/2)\}\le T\rho$
for every $H\in[H_0,H_m]$ and every such law, including $\eta=0$.
\end{theorem}
\begin{proof}
Expand the actual fourth power at each frequency. The linear part is
controlled by the one common measure, and the remaining complex part
by its full nonnegative modulus bound. Midpoint quadrature yields
$G_P(H)-T\rho\le P_H(\eta)-E_H$ at each node.
Convexity gives $P_H(\eta)\le0$ on $[L,U]$. The chord inequality
\eqref{r4bridge:H} now implies $G_P(H)-T\rho\le0$ between nodes.
Finally use \eqref{r4bridge:prawitz}. All these inequalities hold for
one arbitrary law throughout; no law is reconstructed from a relaxation.
\end{proof}

An elementary permitted remainder is obtained from
\[
 \psi^4-g^4-4g^3(\psi-g)
  =(\psi-g)^2(\psi^2+2\psi g+3g^2).
\]
If $|g|\le a\le1$ and $|\psi-g|\le A+B\eta$ nonnegatively on
the whole cell, freeze $p=\min(1,a+\max(A+BL,A+BU))$.
Then $(A+B\eta)^2(p^2+2pa+3a^2)$ is a convex quadratic remainder.
Other complete bounds, including a proven disk/lens bound, may replace
it. Their correctness remains a premise if not proved separately.
Selecting a whole valid row independently at each node is safe;
interpolating a pointwise minimum of candidate curves is not needed.

For $H\le-4$, Cantelli gives $F_Z(H/2)\le1/5$.
For $H\ge25/8$, the elementary normal bound $\Phi(25/16)>4/5$
gives the same raw numerator bound $2/5$. Since $\rho\ge1$, both
positive-score tails are at most $2/5<T=4097/10000$.
To verify that normal bound directly, $\pi<22/7$ gives
$\phi(0)>997/2500$ by squaring. With $x=25/16$, the exponential
Taylor lower bound gives
\[
 \Phi(x)\ge\frac12+\frac{997}{2500}
       \left(x-\frac{x^3}{6}+\frac{x^5}{40}-\frac{x^7}{336}\right)
       >\frac45.
\]

\subsection{The independent count route and reflection}
For this separate count route assign zero to the nonnegative class:
let $p=P(X\ge0)$, $q=1-p$. This differs from the strict-positive
convention of the main argument when there is a zero atom. Both are positive. Conditioning on the
number $j$ of nonnegative observations gives mass $\binom4j p^jq^{4-j}$.
Conditional on any fixed sign pattern, independence retains independent
conditional marginals. Thus $Z$ has conditional mean and variance
\[
 \mu_j=\frac{u(j-4p)}{2pq},\qquad
 V_j=\frac14\left[j\left(\frac{v_+}{p}-\frac{u^2}{p^2}\right)
 +(4-j)\left(\frac{v_-}{q}-\frac{u^2}{q^2}\right)\right].
\]
The same moments hold after mixing patterns with that count. Closed-event
Cantelli yields the source bound $A_j(x)=1$ for $x\ge\mu_j$ and
$A_j(x)=V_j/(V_j+(\mu_j-x)^2)$ otherwise. This includes $V_j=0$.
Cauchy gives
$w(1-\vartheta)/2\le p\le1-w(1+\vartheta)/2$.
Taking upper bounds on all five $A_j$ on a cell and then a valid upper
bound of their single binomial mixture preserves the original-law bound.
No arbitrary dependent sign allocation can replace the iid weights.

Any complete count-based above-cut bound, or independent all-law
large-$\rho$ bound, can join the low-excess theorem when both closed
ranges cover their common cutoff. A parent geometry may be replaced by
all four closed children, each with its own complete low/high cover.
No unexecuted terminal inequality is supplied by this observation.

For the reflected law $\widetilde P=\mathcal L(-X)$, $a,b$ are swapped,
$\vartheta$ changes sign, $\eta$ and $\rho$ are unchanged, and
\[
 \widetilde p=P(X\le0)=1-p+P\{0\},\qquad
 \widetilde\psi(q)=\overline{\psi(q)}.
\]
The raw phase $e^{-\mathrm iqH}$ is retained when evaluating its own
threshold $H$. Applying both positive closed bounds covers the negative
CDF discrepancy because
\[
 \Phi(x)-F_Z(x)\le\Phi(x)-F_Z(x-)
     =F_{-Z}(-x)-\Phi(-x).
\]
Equality at an atom is not assumed. These facts supply the analytic
coverage implication only. Complete numerical enclosures, their selected
implementation, and the unrestricted-owner localization remain separate.
